\section{Theorem and Proof Integration}
The theorem registry is the bridge from named claims to proof sheets, Lean declarations, finite witnesses, and scope boundaries. A theorem row without this bridge is not a promoted theorem.

This section turns selected evidence records into a readable scientific route. The main prose states what is being claimed, what kind of support carries it, and what condition would reopen it. Complete records remain in the evidence package so that the monograph can stay argumentative rather than archival.

\subsection{T133-K0-RES}
The scientific reading is as follows. Claim boundary: K0 resolution-relative distinguishability theorem K0 support is treated as a bounded formal consistency check over declared resolution quotients: same-resolution states are not distinguished, and a finite countermodel shows raw separation need not induce resolution distinction. It is promoted only as a bounded model-core theorem claim, not as an independent novelty or unrestricted theory-wide theorem. Evidence support: public proof sheet for the K0 resolution theorem Lean declaration k0\_countermodel\_raw\_separation\_not\_resolution\_distinction the K0 Resolution Foundation appendix in the master monograph A proof that assumes every pair of raw real states is epsilon-separated is outside current formal profile and fails release criterion. The release consequence is local: this entry supports only the stated claim under its declared assumptions and can be reopened at the named boundary.

\subsection{T133-OMEGA-STATUS}
The evidence interpretation is deliberately bounded. Claim boundary: Typed liveness, death, residue, and rebirth evidence-consistency theorem Any claim reading residue-preserving restart as same-identity survival is rejected unless endpoint-bound identity evidence is supplied. Evidence support: Death blocks live status; residue and rebirth are token-bound evidence relations with distinct class-specific endpoint rules: residue separates source from residue while returning to the source endpoint, and rebirth separates source, residue, and new target tokens. Rebirth is non-identity unless endpoint-bound identity evidence has identity class, declared i. The evidence package records the complete field. public proof sheet for the lifecycle status theorem Lean declaration lifecycle\_residue\_rebirth\_morphism\_boundary content/OC\_1\_3\_3\_TYPED\_FOUNDATION.tex The important distinction is between the public statement, the artifact that carries it, and the condition that would make the statement too strong.

\subsection{T133-K-ZERO}
The reviewer route starts from support, not from assertion. Claim boundary: Continuumness zero obstruction theorem Evidence support: public proof sheet for the K-zero boundary theorem Lean declaration continuumness\_zero\_case\_iff\_declared\_zero\_cause\_with\_support the Continuumness Functionals appendix in the master monograph Reopening boundary: Continuumness zero requires live support, an independently clear obstruction register, and a declared zero-cause family; a zero-cause label alone does not compute k=0. A realization with undeclared zero-cause, missing live support, or any active obstruction may not promote k=0. A hostile reader should test the support class first, then ask whether the wording stays within that support class.

\subsection{T133-BOUNDARY}
The boundary reading prevents overclaim. Claim boundary: Metric-threshold boundary specialization theorem Evidence support: public proof sheet for the boundary representation theorem Lean declaration metric\_boundary\_specialization the Boundary Representation Theorem appendix in the master monograph Reopening boundary: Metric thresholds are a specialization of typed classifier boundaries. A social or logical boundary represented numerically without a measurement rule is blocked by release criterion. This entry is included because it constrains promotion; it does not license a stronger conclusion than the proof, finite case, replay row, or comparator can carry.

Synthesis checkpoint 4. The preceding entries should be read together: promotion is earned by an alignment between claim wording, evidence class, and reopening rule. If any one of those three moves, the public sentence has to move with it.

\subsection{T133-HYBRID}
The scientific reading is as follows. Claim boundary: Typed update and chart-labelled operator semantics theorem Evidence support: OC operators are typed update semantics; chart-labelled flow-one notation is admitted only for declared chart records, while proof/rewrite and guard/reset updates remain first-class non-smooth cases. No differentiability or ODE-solution theorem is promoted in current formal profile. public proof sheet for the hybrid semantics theorem Lean declaration integrated\_operator\_semantics content/OC\_1\_3\_3\_OPERATOR\_SEMANTICS.tex Any route treating a chart token as differentiability, manifold, vector-field, or ODE-solution evidence is outside the current formal profile operator theorem and remains a future lint/proof obligation. The release consequence is local: this entry supports only the stated claim under its declared assumptions and can be reopened at the named boundary.

\subsection{T133-DIM}
The evidence interpretation is deliberately bounded. Claim boundary: Historical axis and effective-rank compatibility theorem Evidence support: Historical axis activation and effective working rank are kept as distinct bounded formal consistency fields; a finite/Lean witness shows compatibility of monotone historical bookkeeping with decreasing effective rank, but current formal profile does not promote an independent scientific dimension theorem. public proof sheet for the dimension semantics theorem Lean declaration historical\_axis\_survives\_rank\_drop the K Level Irreducibility Atlas appendix in the master monograph A K-level is demotable only when the alleged new axis has no witness and no observable consequence. The important distinction is between the public statement, the artifact that carries it, and the condition that would make the statement too strong.

\subsection{T133-CYCLE}
The reviewer route starts from support, not from assertion. Claim boundary: Live-status cycle-mode requirement theorem Evidence support: public proof sheet for the cycle-mode theorem Lean declaration eligible\_live\_requires\_cycle\_or\_maintenance content/OC\_1\_3\_3\_CYCLE\_TAXONOMY.tex An artifact that never updates, replays, checks, or maintains itself is archive residue, not live continuum. Reopening boundary: Declared eligible-live status requires an explicit cycle mode or non-vacuous maintenance predicate. A hostile reader should test the support class first, then ask whether the wording stays within that support class.

\subsection{T133-ID}
The boundary reading prevents overclaim. Claim boundary: Identity, residue, and rebirth evidence-classification theorem Any public claim reading rebirth as literal same-identity survival is blocked. Evidence support: Identity continuation requires endpoint-bound identity evidence: identity class, declared invariant preservation, no residue token, equal source/target endpoint evidence, lifecycle identity-invariant truth, and typed source/target binding. Residue and rebirth evidence classes do not become identity continuation merely by preserving some invariants. No categorical Hom/composition theorem is promoted. public proof sheet for the identity and rebirth theorem Lean declaration endpoint\_bound\_identity\_classification the Identity Residue Rebirth Classification appendix in the master monograph This entry is included because it constrains promotion; it does not license a stronger conclusion than the proof, finite case, replay row, or comparator can carry.

Synthesis checkpoint 8. The preceding entries should be read together: promotion is earned by an alignment between claim wording, evidence class, and reopening rule. If any one of those three moves, the public sentence has to move with it.

\subsection{T133-MIN}
The scientific reading is as follows. Claim boundary: Declared semantic-verdict component independence theorem Evidence support: Within the declared current release tuple semantics, each tuple component has a one-field semantic keep/drop witness that changes the release verdict. public proof sheet for the minimality witness theorem Lean declaration release\_tuple\_semantic\_component\_irredundant the Component witness independence for the declared semantic verdict suite Witnesses appendix in the master monograph Reopening boundary: If any declared tuple component can be removed while `FM-MIN-*` still returns bounded evidence recorded for the declared semantic verdict suite, the minimality card fails. The release consequence is local: this entry supports only the stated claim under its declared assumptions and can be reopened at the named boundary.

\subsection{T133-KLEVEL}
The evidence interpretation is deliberately bounded. Claim boundary: Declared adjacent K-level atlas/evaluator consistency theorem Evidence support: Every declared adjacent K-level transition K0-\textgreater{}K12 has a release-atlas row, retained-witness evaluator check, executable finite row, and inert-witness demotion control inside the current formal profile release classifier; independent semantic irreducibility beyond this declared classifier is a future proof obligation, not a promoted current formal profile theorem. public proof sheet for the K-level witness theorem Lean declaration release\_atlas\_manifest\_has\_total\_finite\_case\_coverage the K Level Irreducibility Atlas appendix in the master monograph If any adjacent K row lacks row identity, retained-witness evaluator failure, or inert-witness demotion control, the finite negative control fails; independent semantic irreducibility remains unpromoted unless separately proven. The important distinction is between the public statement, the artifact that carries it, and the condition that would make the statement too strong.
